\documentclass[11pt,a4paper]{article}
\usepackage[T1]{fontenc}
\usepackage[utf8]{inputenc}
\usepackage{lmodern}
\usepackage[english]{babel}
\usepackage{microtype}
\usepackage[a4paper,margin=2.25cm]{geometry}
\usepackage{xcolor}
\usepackage{enumitem}
\usepackage[most]{tcolorbox}
\usepackage{fancyhdr}
\usepackage{lastpage}
\usepackage{titlesec}
\usepackage{newunicodechar}
\usepackage{hyperref}
\newunicodechar{’}{\textquoteright}
\newunicodechar{“}{\textquotedblleft}
\newunicodechar{”}{\textquotedblright}
\newunicodechar{×}{\ensuremath{\times}}
\definecolor{navy}{HTML}{17365D}
\definecolor{reviewbg}{HTML}{F3F5F7}
\definecolor{reviewframe}{HTML}{AAB4BF}
\definecolor{responsebg}{HTML}{F7FBF8}
\definecolor{responseframe}{HTML}{9CB8A4}
\definecolor{locationbg}{HTML}{FBFAF5}
\definecolor{locationframe}{HTML}{C5B98F}
\definecolor{textgray}{HTML}{4B5563}
\hypersetup{
  colorlinks=true,
  linkcolor=navy,
  urlcolor=navy,
  pdftitle={Response to Reviewers -- Paper 306},
  pdfauthor={Tomasz Piłka, Kaj Skubiszak, Piotr Andrzejewski, Patryk Żyliński}
}
\pagestyle{fancy}
\fancyhf{}
\lhead{Response to Reviewers -- Paper 306}
\rhead{MLSA 2026}
\cfoot{\thepage\ of \pageref{LastPage}}
\renewcommand{\headrulewidth}{0.4pt}
\titleformat{\section}
  {\Large\bfseries\color{navy}}
  {}{0pt}{}
\titleformat{\subsection}
  {\large\bfseries\color{navy}}
  {}{0pt}{}
\setlength{\parindent}{0pt}
\setlength{\parskip}{5pt}
\setlist[itemize]{leftmargin=1.5em,itemsep=2pt,topsep=3pt}
\setlist[enumerate]{leftmargin=1.8em,itemsep=3pt,topsep=3pt}
\newtcolorbox{reviewbox}[1][]{
  enhanced, breakable, colback=reviewbg, colframe=reviewframe,
  boxrule=0.6pt, arc=1.5mm, left=3mm, right=3mm, top=2.5mm, bottom=2.5mm,
  fonttitle=\bfseries, coltitle=navy, title=#1
}
\newtcolorbox{responsebox}{
  enhanced, breakable, colback=responsebg, colframe=responseframe,
  boxrule=0.6pt, arc=1.5mm, left=3mm, right=3mm, top=2.5mm, bottom=2.5mm
}
\newtcolorbox{locationbox}{
  enhanced, breakable, colback=locationbg, colframe=locationframe,
  boxrule=0.6pt, arc=1.5mm, left=3mm, right=3mm, top=2mm, bottom=2mm
}
\newcommand{\response}[2]{%
  \begin{responsebox}
  \textbf{Response:}
  #1
  \end{responsebox}
  \begin{locationbox}
  \textbf{Location in the revised manuscript:}
  #2
  \end{locationbox}
}
\begin{document}
\begin{center}
  {\LARGE\bfseries\color{navy} Response to Reviewers}\\[5pt]
  {\large Paper ID 306}\\[10pt]
  {\Large\bfseries Prediction or Revaluation? Phase- and Space-Aware VAEP for Women's Football}\\[8pt]
    \begin{tabular}{@{}l@{}}
  Tomasz Piłka, Kaj Skubiszak, Piotr Andrzejewski, and Patryk Żyliński
  \end{tabular}
\end{center}
\vspace{4mm}
\hrule
\vspace{5mm}

\section{Summary of changes, and a correction}

We thank both reviewers for reviews that were unusually specific and, as it turned out,
unusually productive. Reviewer~2's Comment~8(b) asked us to state explicitly which action
universe Tables~2 and~3 use. Checking this against our evaluation code, we found that they
do not use the same one, and that the same inconsistency affects the player analysis in
Section~4.3 more seriously still. We describe the error and its consequences here rather
than distributing it across individual responses, because it changes several numbers the
reviewers commented on.

\textbf{The error.} PS-VAEP is defined only on the 51{,}172 actions carrying a 360 freeze
frame, whereas E-VAEP and P-VAEP are defined on all 60{,}370 evaluation actions.
Section~3.6 of the submitted paper states that all ablation metrics are reported on the
matched 51{,}172-action subset, and warns that mixing universes would distort the
comparison. The discrimination rows of Table~3 do follow that rule. Table~2, the ECE row
of Table~3, and all of Section~4.3 do not.

\textbf{Consequences.} Three sets of numbers change.

\begin{itemize}
  \item \textbf{Table~2} now reports all variants on the matched subset. Pooled training
    remains the best strategy for all three variants, but under men-only and fine-tuned
    transfer P-VAEP is no longer above E-VAEP.
  \item \textbf{The ECE row of Table~3} was taken from a calibration table computed on the
    full set. Recomputed correctly, two of its comparisons reverse direction. We no longer
    claim a calibration advantage for any variant, since none of the differences survive
    resampling under either equal-width or adaptive binning.
  \item \textbf{Section~4.3 was materially affected.} Player totals were summed over
    60{,}370 actions for E-VAEP and 51{,}172 for PS-VAEP, so roughly 9{,}200 actions
    contributed to a player's E-VAEP total but not her PS-VAEP total. Recomputed on
    matched actions with the same 117 players, the headline correlation rises from
    $\rho = 0.69$ to $\rho = 0.83$, and the named examples largely dissolve: the rank gains
    we reported as $+102$, $+94$, $+88$ and $+58$ become $+3$, $+21$, $+5$ and $+1$. Those
    figures measured the missing 15\% of each player's actions, not the spatial features.
    We have replaced them with the matched risers, which are positionally mixed rather than
    uniformly centre-forwards.
\end{itemize}

\textbf{What survives.} The paper's central claim is unaffected and is now on firmer
footing. Under the paired match-level bootstrap both reviewers requested, the improvement
of P-VAEP over E-VAEP on conceding ROC-AUC ($+0.023$, 95\% CI $[+0.008, +0.039]$) is the
only difference anywhere in the ablation whose interval excludes zero. Spatial features
leave aggregate prediction unchanged while reordering the player ranking from
$\rho = 0.97$ to $\rho = 0.83$. The prediction-versus-revaluation distinction is therefore
supported by the corrected analysis more sharply than by the submitted one, because the
predictive and revaluation effects now separate cleanly rather than by degrees.

\textbf{Scope.} Between them the reviews propose roughly ten additional experiments. We
have prioritised those that test the paper's claims directly and can be done rigorously
within a workshop revision: the matched recomputation, bootstrap intervals throughout, the
stratified analysis of the revaluation mechanism, and the freeze-frame visibility controls.
We have implemented the phase~$\times$~action-type interactions that Section~3.3 described
but that were absent from our feature code. Where we have not carried out a requested
experiment, we say so plainly below rather than promising it.

\clearpage
\section{Reviewer \#1}

\begin{reviewbox}[Comment 1 --- motivation for phase annotation]
The practical motivation for adding tactical phase annotations could be clearer. The statement that two actions with the same event description may differ in difficulty and value seems more directly related to spatial context. A concrete example where the absence of phase information leads to an implausible valuation would help.
\end{reviewbox}
\response{The reviewer identified a real conflation. The ``same event description, different difficulty'' argument does motivate \emph{spatial} context, and we used it to motivate both layers. The two motivations are now separated.

Phase encodes the \emph{possession regime} a state belongs to, not its local geometry. Canonical VAEP observes a fixed window of three preceding actions, which represents the recent past but not whether the possession is settled or has just changed hands. Two own-third passes can agree on type, body part, start and end coordinates, score state and all three preceding actions, and still carry very different conceding risk depending on whether they occur in settled build-up or three seconds after a turnover with the opponent pressing forward. That is precisely the distinction the phase layer supplies, and it explains why its effect is concentrated on the conceding head --- the one place in our ablation where a robust improvement appears.

We have also stated the applications that require phase annotation rather than merely benefit from it: role-specific player evaluation, where a deep-lying midfielder's build-up contribution should not be pooled with her defensive-transition contribution because the two carry different risk profiles; and phase-specific team and opponent scouting, where the phase is already the standard analytical unit and an unconditioned VAEP total cannot be decomposed into one.}{Section~1, paragraph~2 split into three, with separate motivations for phase and space and the applications stated explicitly.}

\begin{reviewbox}[Comment 2 --- chronology of the pass-valuation citation]
The related-work section states that subsequent work extended action valuation to pass valuation and cites reference [2]. Although reference [2] was published in the same year as the original VAEP paper, the underlying work predates VAEP, so this wording may be misleading.
\end{reviewbox}
\response{We agree; the wording was inaccurate. Bransen et al.\ is not a successor to VAEP, and describing it as ``subsequent work'' misrepresents the chronology. Pass-valuation work is now presented as a parallel line of research on valuing specific action types, developed alongside rather than after the VAEP framework, with later learning-paradigm extensions described separately.}{Section~2.1, final paragraph.}

\begin{reviewbox}[Comment 3 --- aggregate metrics vs.\ specific situations]
The evaluation focuses mainly on aggregate metrics, while the discussion suggests that spatial features matter most in specific high-risk situations. A systematic analysis of the types of actions or situations that benefit most would strengthen the conclusions.
\end{reviewbox}
\response{This was the main evidential gap, and Reviewer~2 raised it as Comment~5. We have added a stratified analysis (new Section~4.4), and it does not support our original claim in full.

We decompose the per-action change $\Delta = V_{\mathrm{PS}}(a) - V_{\mathrm{P}}(a)$ across defensive density, nearest-opponent distance, pitch zone, action type and outcome, reporting the share of actions, the mean change with match-level bootstrap intervals, and within-stratum conceding discrimination.

The result supports a congestion mechanism but not a pressure mechanism. Actions with three or more visible opponents within 5~m --- 5.2\% of the evaluation set --- are the \emph{only} stratum with a robustly positive mean change ($+0.0023$, CI $[+0.0005, +0.0039]$). Distance to the nearest opponent produces no gradient at all ($-0.0015$, $-0.0018$, $-0.0014$, $-0.0013$ across $<2$, $2$--$5$, $5$--$10$, $>10$~m), and neither does pitch zone. By action type, the largest robust negative shifts fall on throw-ins and goalkeeper claims rather than on open-play progression.

We have therefore narrowed the claim in the paper: the redistribution is congestion-driven, and the ``actions under pressure in congested high-value areas'' formulation is replaced by what the strata actually show. The Maanum case study is retained but demoted to an illustration of a mechanism now established quantitatively.}{New Section~4.4 with the stratified table; Section~4.3 and Section~5 rewritten around the narrower claim; case study reframed.}

\begin{reviewbox}[Comment 4 --- feature importance, grouped ablation, event-data approximation]
A feature-importance or grouped ablation analysis would help clarify the contribution of the added phase and spatial features. It would also be useful to assess whether the most informative spatial features could be approximated from event data or another source when 360 data is unavailable.
\end{reviewbox}
\response{We agree that both analyses would be valuable, and we have not carried them out for this revision. We prefer to say so rather than promise them.

Our reasoning is that the grouped ablation requires roughly ten additional training runs and the event-data approximation requires a model per spatial feature plus a rebuilt VAEP variant, and neither would change the paper's conclusion: with the spatial layer as a whole producing no robust aggregate effect, decomposing that effect into groups would mostly resolve noise into finer noise. The bootstrap intervals we now report make this explicit --- the PS-VAEP versus P-VAEP difference on conceding ROC-AUC is $[-0.037, +0.015]$, so a per-group decomposition of it would not be interpretable at this sample size.

The second question we consider the more valuable of the two, precisely because it determines whether the approach transfers to competitions without 360 coverage, and we have flagged it as the primary direction for future work rather than leaving it implicit. We note that our stratified analysis gives a partial answer: the redistribution is driven by local defensive density within 5~m, which is the spatial quantity least likely to be recoverable from event data alone, whereas the coarse zone-context features that events could approximate are not where the effect lies.}{Section~5, future-work paragraph. We have not added the grouped ablation or the approximation experiment.}

\begin{reviewbox}[Comment 5 --- numerical precision and uncertainty]
Table 3 uses inconsistent numerical precision. More importantly, the paper places considerable emphasis on small differences between models. Without uncertainty estimates, some of these differences may not be robust.
\end{reviewbox}
\response{Both points are well taken, and the second turned out to be the more consequential of the two.

\textbf{Precision.} Discrimination and probability metrics are now reported to three decimal places throughout, and ECE in units of $10^{-3}$ so that it is legible at comparable precision. Bolding no longer marks nominal ``best'' values; it marks differences whose interval excludes zero.

\textbf{Uncertainty.} We now report paired match-level bootstrap intervals ($B = 1000$ resamples of the 31 evaluation matches) for every contrast in Tables~2 and~3. Because all variants share an evaluation action set, we resample matches once per replicate and recompute every variant on that resample, so the intervals describe the \emph{difference} between variants, which is what our claims concern.

The reviewer's suspicion was correct. Of every difference in the pooled ablation, exactly two have intervals excluding zero: P-VAEP over E-VAEP on conceding ROC-AUC ($+0.023$, CI $[+0.008, +0.039]$), and PS-VAEP's \emph{degradation} of scoring ROC-AUC ($-0.005$, CI $[-0.010, -0.0002]$). The conceding-AP advantage of PS-VAEP that we emphasised is not robust ($+0.006$, CI $[-0.027, +0.038]$), nor is any log-loss or calibration difference. In Table~2, only the pooled P-versus-E difference is robust.

We now present this as a finding rather than a caveat. That phase produces the ablation's only robust predictive gain while space produces none, yet space alone reorders the player ranking, states the paper's prediction-versus-revaluation thesis more sharply than the submitted version managed.}{Section~3.6, bootstrap protocol; Tables~2 and~3 reformatted with uniform precision and paired intervals; Sections~4.1--4.3 and~5 rewritten to distinguish robust from non-robust differences.}

\begin{reviewbox}[Comment 6 --- construction of the men's source corpus]
The construction of the men's source dataset could be motivated better. It combines men's international tournaments with a domestic Bundesliga season. Competition type may be at least as important as gender for domain similarity.
\end{reviewbox}
\response{We accept the criticism. The source corpus was assembled to maximise available men's 360 data --- Bundesliga 2023/24 is the largest men's 360 corpus in the public StatsBomb release, and excluding it would have reduced the source domain substantially --- and the submitted paper presented that pragmatic choice as though it were a design decision. Section~3.2 now states the actual reason.

On the substantive hypothesis, we have not run the international-versus-club transfer comparison or estimated domain divergence, and we do not want to claim otherwise. We agree the hypothesis is testable with data we already hold and that it is a natural next step; we have added it to future work. We would rather report it properly than attach an underpowered version of it to a revision in which the transfer comparisons already carry intervals spanning zero.}{Section~3.2, corpus motivation; Section~5, future work. The competition-type comparison has not been run.}

\begin{reviewbox}[Comment 7 --- relationship to VAEP360 / Un-XPass]
The relationship to the VAEP360 approach used in the Un-XPass paper should be discussed more explicitly. A direct comparison with Phase-Space VAEP would help position the contribution.
\end{reviewbox}
\response{We agree the positioning was too implicit. Section~2.2 now discusses the relationship along three axes: \emph{objective} --- VAEP360 as used in Un-XPass supports evaluation of passing decisions and creativity, whereas our target is the full on-ball action stream including carries, take-ons and defensive actions; \emph{representation} --- VAEP360 consumes the freeze frame through learned surface-based representations, while we use seventeen handcrafted geometric features chosen for interpretability and robustness under partial visibility in a small target corpus; and \emph{data regime} --- our setting is deliberately low-capacity because the target corpus is small and cross-domain.

On a direct empirical comparison we would prefer to be candid: a faithful reimplementation of VAEP360 within our source-target protocol is beyond what we can complete for this revision, and an unfaithful one would be worse than none. We provide a structured conceptual comparison and state the empirical comparison as future work rather than claiming a superiority we have not demonstrated.}{Section~2.2, comparison paragraph; Section~5, final limitation.}

\begin{reviewbox}[Comment 8 --- choice of LightGBM and hyperparameter sensitivity]
The choice of LightGBM and the selected hyperparameters require more justification. Given the small differences between some models, sensitivity to these choices may matter.
\end{reviewbox}
\response{The choice was under-justified, and Section~3.4 now states the reasoning: the feature set is tabular, low-dimensional and mixes categorical with continuous variables; the positive classes are rare ($\approx 1.5\%$ and $0.34\%$); the target corpus is small; and gradient-boosted trees are the standard learner in the VAEP literature and in the reference \texttt{socceraction} implementation, which keeps our results comparable with prior work. We also now state explicitly that hyperparameters were taken from that reference configuration rather than tuned per variant --- a fairness property of the ablation, since all variants share the configuration.

We have not run the multi-seed and hyperparameter-grid sensitivity analysis. We note that the bootstrap intervals now reported address the reviewer's underlying concern from a different direction: they show directly that all but one of the differences we discuss are indistinguishable from zero under resampling of the evaluation data. Seed and hyperparameter variation would add further sources of variability on top of that, and could only widen those intervals, not narrow them. The one difference that survives --- $+0.023$ with a lower bound of $+0.008$ --- is the one for which a sensitivity check would be most informative, and we agree it should be run before any stronger claim is made about its magnitude.}{Section~3.4, extended justification. The sensitivity analysis has not been run.}

\subsection*{Questions}

\begin{reviewbox}[Question 1]
Which concrete applications of VAEP require explicit tactical phase annotations, and can the authors provide an example where missing phase information leads to an implausible valuation?
\end{reviewbox}
\response{Please see our response to Comment~1, which sets out the applications --- role-specific player evaluation, phase-specific scouting, and risk-aware valuation where the cost of a turnover depends on the possession regime --- and the mechanism by which two event-identical own-third passes receive the same valuation despite differing conceding risk.}{Section~1.}

\begin{reviewbox}[Question 2]
In the transfer-learning experiments, is the main benefit simply the availability of more 360 data, regardless of whether it comes from men's or women's football? How is the effect of sample size separated from the effect of cross-domain transfer?
\end{reviewbox}
\response{It is not separated, and on the evidence we have we cannot separate it. Pooled training sees more data than either single-domain strategy, so the pooled advantage is consistent with both cross-domain complementarity and a pure sample-size effect. Reviewer~2 raised the same confound as Comment~3.

Rather than defend the stronger reading, we have narrowed the claim. The paper now states that pooled training outperforms the alternatives and that this does not by itself establish cross-domain complementarity, and identifies the size-matched comparison and per-domain learning curves as the experiments that would settle it. We have not run them for this revision.

We would add that the corrected Table~2 makes the transfer claims weaker than they appeared: on the matched action set, every pairwise difference under men-only, women-only and fine-tuned training has a bootstrap interval spanning zero. Only the pooled P-versus-E difference is robust. The transfer contribution is therefore best read as establishing that women-only training is insufficient, which the data do support, rather than as a ranking among the other three strategies.}{Section~4.1 and Section~5, claim narrowed; Abstract adjusted. Size-matched comparison not run.}

\begin{reviewbox}[Question 3]
Which actions or tactical situations benefit most from the phase and spatial features? Can the authors provide a systematic analysis of the ``rare high-risk cases''?
\end{reviewbox}
\response{Please see our response to Comment~3 for the stratified analysis. On the specific phrase ``rare high-risk cases'': we withdraw it. We used it to gloss PS-VAEP's higher conceding AP, and on inspection that AP advantage rests on a handful of actions at the very top of the ranking --- PS-VAEP retrieves 8 of 175 concessions in its top 10 ranked actions against P-VAEP's 6 --- while being no better than P-VAEP at every threshold from the top 1\% to the top 50\%. That is not a systematic advantage on rare high-risk cases, and its bootstrap interval spans zero. The paper now says so.}{New Section~4.4; Section~4.2 rewritten.}

\begin{reviewbox}[Question 4]
What is the individual contribution of the different phase and spatial feature groups? Could some of the most informative spatial features be approximated from event data?
\end{reviewbox}
\response{Please see our response to Comment~4. We have not run the grouped ablation or the approximation experiment, and we explain there why we judged them uninformative at this sample size and what partial answer the stratified analysis provides.}{Section~5, future work.}

\begin{reviewbox}[Question 5]
Why was Bundesliga 2023/24 included alongside men's international tournaments? Did the authors examine whether competition type is more important than gender for source-target similarity?
\end{reviewbox}
\response{Bundesliga 2023/24 was included because it is the largest men's 360 corpus in the public StatsBomb release; the decision was driven by data availability and the submitted paper should have said so. We had not examined whether competition type matters more than gender, and we have not examined it for this revision either. See our response to Comment~6.}{Section~3.2; Section~5, future work.}

\begin{reviewbox}[Question 6]
How does Phase-Space VAEP differ from the VAEP360 approach used in Un-XPass? Could VAEP360 also be extended with the proposed phase annotations?
\end{reviewbox}
\response{Please see our response to Comment~7 for the full comparison. On the second question --- yes. The phase layer is a property of the game state, not of the spatial encoder, and is therefore orthogonal to how the freeze frame is represented; it could be supplied to VAEP360 as additional conditioning input without modification. Our results suggest this is worth testing, since in our setting the phase layer contributed the only robust predictive gain of the two.}{Section~2.2; Section~5.}

\begin{reviewbox}[Question 7]
Why was LightGBM selected, and how sensitive are the results to the chosen hyperparameters and random variation?
\end{reviewbox}
\response{Please see our response to Comment~8 for the selection rationale, now stated in Section~3.4, and for our reasoning about sensitivity, which we have not measured directly but which the bootstrap intervals bound from one direction.}{Section~3.4.}

\clearpage
\section{Reviewer \#2}

\begin{reviewbox}[Comment 1 --- complete the ablation with Event+Space]
Complete the ablation. The nested comparison omits an Event+Space model. The current design measures the marginal effect of space after phase has already been added, but cannot separate the main effects of phase and space from their interaction. A 2×2 Event / Phase / Space ablation would support the paper's central argument more directly.
\end{reviewbox}
\response{The reviewer is right that the nested design confounds the main effect of space with its interaction with phase, and that this is a design flaw rather than a reporting omission. The nested design was chosen because phase is far cheaper to compute than freeze-frame features, but that is a reason about cost, not about inference.

We have not trained the Event+Space variant for this revision. Our reason is that the bootstrap analysis changes what the missing cell could tell us: since the marginal effect of space after phase is $[-0.037, +0.015]$ on conceding ROC-AUC and $[-0.027, +0.038]$ on conceding AP, decomposing an effect of that size into main effects and an interaction would not be identifiable at $n = 31$ matches. We would be reporting three numbers whose intervals all contain zero.

We think the reviewer's underlying point is nonetheless addressed by the revision, though by a different route. The prediction-versus-revaluation distinction predicts that the phase effect should appear in probability metrics and the space effect in player-ranking correlations, and the corrected analysis shows exactly that separation: phase gives the ablation's only robust predictive gain and barely moves the ranking ($\rho = 0.97$), while space gives no robust predictive gain and moves the ranking substantially ($\rho = 0.83$). We would welcome guidance on whether the workshop would prefer the Event+Space cell reported with intervals spanning zero.}{Section~4.2 and Section~5 rewritten around this separation. ES-VAEP not trained.}

\begin{reviewbox}[Comment 2 --- quantify uncertainty]
Quantify uncertainty. Several conclusions rely on small metric differences and one held-out tournament. Please report paired match-level bootstrap confidence intervals for Tables 2 and 3. For player rankings, resample matches and recompute player aggregates and correlations rather than treating players as independent observations.
\end{reviewbox}
\response{We have adopted exactly the protocol specified, and it changed our conclusions. Reviewer~1 raised the same concern as Comment~5, where we report the ablation results.

For Tables~2 and~3 we resample the 31 evaluation matches with replacement ($B = 1000$) and recompute all variants on each resample, forming paired differences within each replicate. Two differences out of the whole ablation have intervals excluding zero, both reported in our response to Reviewer~1's Comment~5.

For the player analysis we follow the prescription exactly: matches, not players, are the resampling unit, and within each replicate we recompute per-player minutes and VAEP per~90 before recomputing the Spearman correlations and individual rank changes. The reviewer's anticipation was correct --- the intervals are wide. Of the 117 qualifying players, only 32 have a rank change whose interval excludes zero. We report that number in the paper and state the positional conclusions only to the extent the intervals support them.

We should add that this analysis is what led us to find the action-universe error described on page~1, since the matched correlations came out substantially higher than the ones we had reported.}{Section~3.6, protocol; Tables~2, 3 and~4 with intervals; Section~4.3 with rank-change intervals.}

\begin{reviewbox}[Comment 3 --- strengthen or narrow the transfer claim]
Strengthen or narrow the transfer-learning claim. Pooled training is evaluated only through conceding ROC-AUC and also uses more total training data than either single-domain strategy. The current result therefore does not distinguish cross-domain complementarity from a sample-size effect. Reporting both heads and probability metrics, alongside a size-matched or domain-balanced comparison, would make this contribution more convincing.
\end{reviewbox}
\response{We accept both parts and have narrowed rather than defended the claim. Reviewer~1's Question~2 raises the same confound, and our fuller answer is there.

On reporting: selecting a single metric for a single head was a reporting decision we cannot justify, and both heads with AP, Brier, log loss and ECE are now available for all four strategies and all variants, with bootstrap intervals.

On the confound: we have not run the size-matched or domain-balanced comparison, and we therefore do not claim cross-domain complementarity. The paper now claims only what the data support --- that women-only training is insufficient for this target domain --- and identifies the size-matched comparison as the experiment needed to go further.}{Section~4.1; claim narrowed in Abstract, Section~1 and Section~5. Size-matched comparison not run.}

\begin{reviewbox}[Comment 4 --- clarify the spatial prediction result]
Clarify the spatial prediction result. PS-VAEP improves conceding AP monotonically, despite worsening ROC-AUC and log loss relative to P-VAEP. This suggests a potentially interesting trade-off in ranking rare positive cases rather than a general predictive improvement. The paper should highlight this distinction and quantify its uncertainty. Given the very low conceding rate, reliability diagrams or adaptive-bin calibration metrics would also support the calibration claims.
\end{reviewbox}
\response{We investigated the reviewer's hypothesis directly, and it is not what is happening --- the pattern is thinner than a ranking trade-off.

Decomposing the conceding ranking, PS-VAEP retrieves 8 of the 175 concessions in its top 10 ranked actions against P-VAEP's 6, but is equal or worse at every threshold from the top 1\% to the top 50\% (at the top 1\%: 0.104 precision against 0.117; at the top 10\%: 0.192 against 0.197). Because average precision weights early recall heavily, a difference of two concessions among the first ten predictions is enough to move it. The AP advantage is an artefact of the sharp end of the ranking rather than a systematic improvement on rare positives, and its bootstrap interval spans zero ($+0.006$, CI $[-0.027, +0.038]$). We now report this decomposition and have withdrawn the ``rare high-risk cases'' framing.

On calibration, we accept that 15 equal-width bins are close to uninformative at a $0.34\%$ positive rate. We now report adaptive (equal-frequency) binning alongside equal-width. The two disagree about which variant is best calibrated --- equal-width favours P-VAEP on the conceding head, adaptive favours E-VAEP on both --- and no calibration difference is robust under resampling. We therefore make no calibration claim for any variant, which also retracts the ``best conceding-head ECE'' statement from the submitted Section~4.2.}{Section~3.6, revised calibration metrics; Section~4.2, top-$k$ decomposition and retracted calibration claim.}

\begin{reviewbox}[Comment 5 --- systematic evidence for the revaluation mechanism]
Provide systematic evidence for the revaluation mechanism. The positional shifts and Maanum example are suggestive, but one case study does not establish that pressure and congestion drive the redistribution. Analyze the change from P-VAEP to PS-VAEP across nearest-opponent distance, defensive density, action type, pitch zone, and outcome. Feature importance could supplement this analysis but would not replace it.
\end{reviewbox}
\response{We agree, and this is the change we consider most important in the revision. The new Section~4.4 analyses $\Delta = V_{\mathrm{PS}}(a) - V_{\mathrm{P}}(a)$ across exactly the five dimensions listed, with bootstrap intervals and within-stratum discrimination. Our full summary of the results is in the response to Reviewer~1's Comment~3.

The short version is that the analysis supports half of our claim and refutes the other half. Congestion is supported: actions with three or more opponents within 5~m are the only stratum with a robustly positive mean change. Pressure is not: nearest-opponent distance produces no gradient. Pitch zone produces none either, so ``congested high-value areas'' is not right --- the effect is about local density, not location. And the largest robust negative shifts are on throw-ins and goalkeeper claims, not on the high-volume open-space progression we described. The paper's claim has been narrowed accordingly.

The positional-shift evidence also weakened once computed on matched actions, for the reasons on page~1: the centre-forward pattern we reported was substantially an artefact of the action-universe mismatch. The Maanum action is retained as an illustration of the congestion stratum, explicitly labelled as illustrating rather than establishing.}{New Section~4.4; Section~4.3 rewritten; case study reframed.}

\begin{reviewbox}[Comment 6 --- partial freeze-frame visibility]
Address partial freeze-frame visibility. Counts of nearby or goal-side opponents depend on which players are visible in the broadcast frame. Since visibility varies with pitch position and camera framing, some positional redistribution may reflect coverage rather than defensive structure. Controlling for visible-area size, number of visible players, or proximity to the visible-area boundary would strengthen the interpretation.
\end{reviewbox}
\response{This was a serious threat to our interpretation and we are grateful it was raised precisely --- broadcast coverage varies systematically with pitch position, so the confound is aligned with our headline result rather than orthogonal to it. We tested it, and the result is favourable.

Coverage does vary by zone in the direction the reviewer anticipates: a mean of 14.4 visible players in the middle third against 12.2 in the penalty area. But the association with the revaluation is negligible. The Spearman correlation between the number of visible players and $\Delta$ is $0.019$ across the evaluation set and $0.065$ within the penalty area, the zone where the confound would bite hardest. Coverage therefore does not plausibly drive the redistribution.

We report the coverage descriptives by zone and the association with $\Delta$ so that readers can judge this directly, and we retain the point as a limitation on what freeze-frame features can encode in general, distinct from the specific confound the reviewer identified.}{New material in Section~4.4; Section~5 limitations updated.}

\begin{reviewbox}[Comment 7 --- interpret redistribution neutrally]
Interpret redistribution neutrally. Reduced credit for progression in open space should not yet be framed as a correction. Such progression may be legitimately valuable, including because preceding off-ball actions created the space. Independent expert validation is necessary before concluding that PS-VAEP allocates value more accurately.
\end{reviewbox}
\response{We accept this and have removed the evaluative framing throughout. The reviewer's point is also internally consistent with a limitation we stated but failed to apply: our model values on-ball actions only, so where the space for a progressive pass was created by preceding off-ball movement, PS-VAEP's reduced credit reflects the model's blind spot rather than a correction of the event-only model. We now make that connection explicitly rather than leaving the limitation isolated.

Concretely: we have removed the phrasing suggesting PS-VAEP allocates value more appropriately; replaced ``tactically plausible redistribution'' with neutral wording describing what changes and in which direction; and stated in both the Abstract and the Discussion that the direction of the redistribution is not evidence of its correctness.

We would note that the stratified analysis has made this easier to honour, since the open-space-progression story it was attached to did not survive (see Comment~5). What we retain is the descriptive claim that credit moves towards actions in locally congested situations, with no normative reading attached.}{Abstract; Section~1, contribution~3; Section~4.3; Section~5, paragraphs~3--4.}

\begin{reviewbox}[Comment 8 --- minor clarity and reporting issues]
\begin{enumerate}[label=\alph*)]
  \item Section 5 contains a duplicated near-verbatim closing passage and the typo ``probability stimation.'' Please tighten this section.
  \item Clarify the action universe used in Tables 2 and 3. The E-VAEP pooled values suggest that both use the matched 51,172-action subset, but this should be stated explicitly.
  \item The n=117 Spearman analysis is descriptive; match-level bootstrap confidence intervals for the correlations and individual rank changes would help assess the stability of the positional conclusions.
  \item An appendix enumerating the 17 spatial features with exact definitions would aid reproduction.
\end{enumerate}
\end{reviewbox}
\response{\textbf{(a)} The duplication was an editing error and the typo was real. Section~5 has been rewritten.

\textbf{(b)} This question is what uncovered the error described on page~1, and we are grateful for it. The inference was reasonable but incorrect: the agreement of the E-VAEP pooled conceding ROC-AUC between the two tables at three decimal places is a coincidence. On the full set it is $0.8374$ and on the matched subset $0.8372$; both round to $0.837$. P-VAEP does not agree between universes ($0.858$ against $0.860$), which is what alerted us. Table~2 as submitted used the full 60{,}370-action set for E-VAEP and P-VAEP and the 51{,}172-action subset for PS-VAEP. All tables in the revision use the matched subset, and both captions and Section~3.6 now state so explicitly.

\textbf{(c)} Addressed by the resampling protocol described under Comment~2: correlations and individual rank changes now carry match-level bootstrap intervals, and only 32 of 117 rank changes have intervals excluding zero.

\textbf{(d)} An appendix enumerating all seventeen spatial features with exact definitions, units, the normalisation convention and the treatment of partially visible players is included if the page allowance permits; otherwise it will be released with the code, which is public.}{Section~5 rewritten; Table captions and Section~3.6 clarified; Table~4 and Section~4.3 with intervals; Appendix~A.}

\vspace{4mm}
\hrule
\vspace{3mm}
We are grateful to both reviewers. The requests for uncertainty quantification and for an
explicit statement of the action universe did more than improve the presentation: together
they surfaced an error that had inflated our headline result, and the corrected analysis
supports the paper's central distinction more cleanly than the submitted one did.

\end{document}
